\chapter{Discussion and conclusion}
\label{cha:discussion-conclusion}

On the circle \(S^1\), functions can be decomposed into sine and cosine waves, i.e. the Fourier basis.
This makes the circle a useful setting for studying spectral bias. Low
frequencies describe slowly varying structure, while high frequencies describe
more oscillatory structure. The question is whether the training dynamics treat
these components differently, and whether this distinction remains meaningful
after moving away from the ideal infinite-width and infinite-data setting.

In the continuum infinite-width reference model, the answer is explicit. The
limiting NTK is rotation-invariant, so the associated operator is diagonal in
the Fourier basis. Each Fourier component of the residual evolves independently:
\begin{equation}
	\alpha_p(t)
	=
	\exp(-\lambda_p t)\alpha_p(0),
	\label{eq:discussion-continuum-mode-decay}
\end{equation}
where \(\alpha_p(t)\) is the residual coefficient of the \(p\)-th Fourier basis
function, and \(\lambda_p\) is the corresponding kernel eigenvalue. For the
ReLU NTK studied here, the low-frequency eigenvalues are larger, so the
low-frequency residual components decay faster. This gives the fixed-kernel
spectral-bias mechanism in this setting
\citep{jacot2018ntk,lee2020wide,rahaman2019spectralbiasneuralnetworks,
	basri2019convergence,bietti2019inductivebiasneuraltangent,
	cao2020understandingspectralbiasdeep}.

The continuum calculation also shows the importance of including a trainable hidden bias in the
network. With the hidden-bias contribution, odd frequencies \(k\ge 3\) have
nonzero eigenvalues. Without this contribution, those frequencies are null
directions of the fixed-kernel dynamics. Thus, the architecture itself
affects whether some frequencies are learned at all by the fixed-kernel
operator.

\section{RQ1: Finite sampling}
\label{sec:discussion-finite-sampling}

The first research question asked what happens when the uniform continuum
measure on \(S^1\) is replaced by finitely many sampled input points.

Exact Fourier-wise learning is lost after sampling. The sampled sine and cosine
vectors are not exactly orthonormal, and they are not exact eigenvectors of the
sampled NTK operator. However, on any fixed low-frequency block
\begin{equation}
	\mathcal H_K
	=
	\operatorname{span}\{1,\cos(\varphi),\sin(\varphi),\ldots,
	\cos(K\varphi),\sin(K\varphi)\},
	\label{eq:discussion-low-frequency-block}
\end{equation}
this error decreases as the number of samples increases. At fixed confidence,
the typical finite-sample error is of order \(n^{-1/2}\), up to logarithmic
factors.

This means that the Fourier picture survives as an approximation, not as an
exact statement. The finite-sample corrected eigenvalues describe the dominant
early-time dynamics of the retained Fourier components. Later in training, some
components have already become very small. At that point, the remaining
finite-sample error can become visible relative to the component being plotted.
This explains why the empirical curves can deviate from the ideal prediction at
late times without contradicting the finite-sample theory.

The preconditioning experiment supports the same interpretation. If the
different decay rates are mainly caused by the Fourier eigenvalues, then
rescaling by those eigenvalues should remove much of the separation between the
retained components. This is what is observed. The experiment therefore acts as
a check that the rate separation is mostly explained by the predicted spectrum,
rather than by an unrelated numerical artifact.

The answer to the first research question is therefore that finite sampling does
not preserve exact Fourier diagonalization, but it preserves it approximately on
fixed low-frequency blocks. The approximation improves with sample size, and it
is strong enough to explain the leading residual dynamics seen in the
experiments.

\section{RQ2: Finite width}
\label{sec:discussion-finite-width}

The second research question asked what happens when the network has finite
width. There are two distinct cases. In the frozen-kernel case, the tangent
kernel is sampled at initialization and then kept fixed. In the evolving-kernel
case, the tangent kernel changes during training.

For the frozen tangent kernel, finite width behaves like a random perturbation
of the infinite-width operator. On the fixed low-frequency space
\(\mathcal H_K\), the projected finite-width operator concentrates around the
continuum operator at rate \(m^{-1/2}\), up to logarithmic factors. The
eigenspace experiments show the same effect geometrically: as width increases,
the finite-width eigenspaces become better aligned with the corresponding
Fourier frequency subspaces.

The approximation is not equally stable for every frequency. Low frequencies
have larger eigenvalues, so a small perturbation has less relative effect on
their dynamics. Higher frequencies have smaller eigenvalues and smaller gaps
between nearby eigenvalues, so the same perturbation can be more visible. This
matches the broader spectral-bias picture: high-frequency components are learned
later and are more sensitive to details of the model and data distribution
\citep{rahaman2019spectralbiasneuralnetworks,basri2019convergence,
	basri2020frequencybiasneuralnetworks}.

The evolving-kernel case is different. Here the fixed-kernel explanation is no
longer complete, because the operator itself changes during training. The
results in this part are empirical rather than theorem-level statements.

At small widths, allowing the kernel to evolve improves the final loss. At
larger widths, this advantage becomes smaller and eventually disappears. The
measurements suggest that the improvement is not caused by better alignment with
the Fourier subspaces. Several higher-frequency subspaces become less aligned
during training. Instead, the main visible change is that the evolving kernel
increases the average operator strength in the lowest Fourier blocks, especially
\(k=0\) and \(k=1\), with a smaller increase for \(k=2\). This agrees with the
residual plots, where the strongest improvements occur in low and intermediate
target components.

The amplitude-variant experiments make this interpretation less dependent on
the particular target used in the main experiment. The same low-frequency
strength increase appears when all active non-constant modes have equal
amplitudes, and also when larger amplitudes are assigned to higher frequencies.
This suggests that the effect is not simply caused by the original target having
larger low-frequency coefficients.

The answer to the second research question is therefore split. For frozen finite
width, the continuum Fourier picture remains approximately valid when the width
is large. For evolving finite width, the Fourier subspaces remain useful as
diagnostics, but the fixed-kernel picture no longer fully explains the dynamics.
Training can reshape the tangent kernel, and this can change the effective
learning rates of different frequency components.

\section{Main conclusion}
\label{sec:discussion-main-conclusion}

The main research question asked whether the continuum infinite-width Fourier
description survives under finite sampling and finite width.

For finite sampling, it survives approximately on fixed low-frequency blocks.
The sampled operator is not exactly diagonal in the Fourier basis, but the error
decreases with the number of samples, and the corrected eigenvalues explain the
dominant early-time residual dynamics.

For frozen finite width, it also survives approximately. The finite-width
operator is random at initialization, but its projection onto \(\mathcal H_K\)
concentrates around the continuum operator as the width increases.

For evolving finite width, the Fourier description changes role. It is no
longer a complete dynamical explanation, because the kernel is not fixed. It is
instead a useful coordinate system for measuring how training changes the
operator. In the experiments here, the most important change is an increase in
low-frequency operator strength at small width.

Thus, finite sampling and frozen finite width behave like controlled
perturbations of the ideal NTK model on fixed low-frequency blocks. Kernel
evolution adds a separate effect: training can modify the operator itself, and
this can make the finite-width network behave differently from its frozen
kernel approximation.

\section{Limitations}
\label{sec:discussion-limitations}

The concentration results are fixed-block results. They control a prescribed
low-frequency space \(\mathcal H_K\) with fixed \(K\). They do not give a
uniform statement over all Fourier modes, and they are not designed to remain
sharp when \(K\) grows with \(n\) or \(m\). This matters because higher
frequencies have smaller eigenvalues and narrower eigengaps. Controlling more
frequencies would require more accurate operator approximations.

The bounds are also conservative. The proofs use entrywise concentration and
union bounds over the retained block. These arguments are enough to prove the
correct decreasing trends, but they discard matrix structure and variance
information. This explains why the theoretical bounds are much larger than the
observed numerical errors. Sharper estimates would likely require direct
operator-norm bounds, variance-sensitive concentration, or tools for U- and
V-statistics.

The finite-sample theory explains the leading behaviour of the sampled
fixed-kernel dynamics, but it does not give a sharp time-uniform description of
every residual component. Once a component has decayed close to zero, small
operator errors can dominate the plotted quantity. The current bounds explain
why such deviations should occur, but they do not predict the exact level of
each late-time deviation.

The finite-width theorem controls the frozen tangent kernel at initialization.
It does not control the evolving tangent kernel during training. The
evolving-kernel experiments show a consistent pattern, but they do not establish
a theorem for the time-dependent operator. A rigorous explanation would need to
control how the tangent kernel moves and how this movement interacts with the
current residual.

The alignment diagnostics also have a limited interpretation. Projector errors
and principal angles measure how close empirical eigenspaces are to Fourier
subspaces, but they do not by themselves determine the loss dynamics. Loss decay
depends on how the current operator acts on the current residual. This is why
worse Fourier alignment can coexist with lower loss when the evolving kernel
also increases operator strength in important directions.

Finally, the thesis does not prove a joint theorem combining finite sampling,
finite width, and training time. The finite-sample and frozen finite-width
analyses isolate separate perturbations of the continuum operator. Actual finite
training runs contain sampling error, finite-width randomness, and kernel
evolution at the same time.

\section{Beyond the uniform circle setting}
\label{sec:discussion-beyond-uniform-circle}

The Fourier basis is the right reference in this thesis because the domain is
the circle and the input distribution is uniform. These assumptions make the
limiting kernel rotation-invariant, so the associated operator is diagonalized
by sine and cosine functions.

For non-uniform input densities, the Fourier basis should not be expected to
remain the exact reference. The relevant operator is weighted by the input
density, and its eigenfunctions need not be global Fourier modes. In that case,
spectral bias should be understood more generally as a preference for the
leading eigenfunctions of the kernel operator induced by the architecture and
the data distribution
\citep{basri2020frequencybiasneuralnetworks,
	bowman2022spectralbiasoutsidetraining}. These eigenfunctions may correspond
to smooth variation along the data distribution, rather than low frequency in
the usual Fourier sense.

For other domains, the same idea applies. On higher-dimensional spheres, the
analogues of Fourier modes are spherical harmonics. On less symmetric domains,
there may be no closed-form eigensystem. The operator viewpoint still applies,
but the preferred directions may be much harder to describe explicitly.

Changing the activation function would also change the kernel spectrum. With an
isotropic initialization on the uniform circle, activations such as sigmoid or
tanh may still produce rotation-invariant kernels, so the Fourier basis may
remain an eigenbasis. However, the eigenvalues would change. A smoother
activation can produce a smoother kernel, which may make high-frequency
eigenvalues decay faster. Polynomial activations can behave differently again,
because they may only activate finitely many harmonic components. Thus spectral
bias is not independent of the activation. It depends on the spectrum of the
kernel generated by the architecture, parameterization, and data distribution.

The role of bias terms may also change. In the ReLU model studied here, the
trainable hidden bias supplies nonzero eigenvalues to odd modes \(k\ge 3\).
For another activation or parameterization, the same odd-even structure need not
hold. The broader point is that architecture can affect not only the speed of
learning different components, but also which components are visible to the
fixed-kernel dynamics.

Moving away from the NTK parameterization would change the interpretation even
more. The fixed-kernel limit is tied to the NTK scaling. In feature-learning or
mean-field regimes, the representation can move substantially during training,
so a fixed initial kernel cannot be expected to describe the full dynamics
\citep{Mei2018meanfieldview,chizat2018globalconvergencegradientdescent}. In
such regimes, the natural object is no longer only the spectrum at
initialization, but the evolution of the operator or of the parameter
distribution during training.

\section{Future work}
\label{sec:discussion-future-work}

The most direct next step is to develop theory for the evolving finite-width
tangent kernel. The experiments in this thesis compare frozen and evolving
dynamics, but the evolving-kernel analysis remains empirical. A natural goal is
to control the time-dependent low-frequency block, its drift from
initialization, and its effect on the residual dynamics.

A second direction is to prove a joint result for finite sampling, finite width,
and training time. This would better match actual training runs, where all three
effects occur together.

A third direction is to sharpen the concentration bounds. The current proofs are
sufficient to show decreasing error with \(n\) and \(m\), but they are not
numerically sharp. Direct operator-norm concentration, variance-sensitive
estimates, or concentration inequalities for U- and V-statistics may give bounds
closer to the observed errors.

A fourth direction is to study the alignment diagnostics more directly. The
projector-error and principal-angle plots suggest that different Fourier
subspaces have different finite-width stability. A more detailed analysis could
measure which Fourier blocks are most strongly mixed by the finite-width
operator and how this depends on eigenvalues and eigengaps.

Finally, the analysis should be extended beyond the uniform circle. Non-uniform
densities, higher-dimensional domains, other activations, and different
parameterizations all change the relevant kernel operator and its spectrum.
Studying these cases would clarify which conclusions are special to the explicit
Fourier model on \(S^1\), and which reflect a broader mechanism behind spectral
bias.

\section{Conclusion}
\label{sec:discussion-final-conclusion}

This thesis studied spectral bias in a setting where the ideal training dynamics
can be written explicitly. On the uniform circle, the infinite-width NTK is
diagonal in the Fourier basis. Each frequency has its own learning rate, and low
frequencies are learned faster because they correspond to larger kernel
eigenvalues.

The results show that this picture is stable under finite sampling and frozen
finite width, as long as one works on a fixed low-frequency block. In both
cases, the ideal Fourier operator is perturbed, but the perturbation decreases
with sample size or width. The evolving finite-width experiments show where the
fixed-kernel explanation stops being complete. When the kernel moves during
training, the Fourier basis remains useful for measurement, but the dynamics are
also shaped by how the operator itself changes.

The final picture is therefore simple. The Fourier basis is exact in the
continuum reference model. It remains approximately valid for finite sampling
and frozen finite width. For evolving finite width, it becomes a diagnostic for
understanding how training reshapes the kernel. Turning this last observation
into a theorem is the main open problem left by this work.
